(1) 非旋转无电荷黑洞的Hawking温度 \(T_H=\hbar c^3/(8\pi GMk_B)\)。用热力学第一定律 \(d(Mc^2)=T_HdS\) 推导熵 \(S(M)\)，选择 \(S(0)=0\)，并与视界面积 \(A_H=16\pi G^2M^2/c^4\) 联系。
(2) 求定体参数热容 \(C=d(Mc^2)/dT_H\) 的符号与质量标度，解释黑洞向冷环境辐射时温度为何反而上升，以及不能直接套用普通正热容物体稳定性的原因。
(3) 将黑洞近似为灰度因子常数 \(\epsilon>0\) 的Stefan辐射体 \(P=\epsilon\sigma A_HT_H^4\)，其中 \(\sigma=\pi^2k_B^4/(60\hbar^3c^2)\)。推导 \(P(M)\) 与质量损失方程，严格求从 \(M_0\) 到 \(M\) 的时间。
(4) 在此近似下求蒸发总时间及 \(T_H(t)\) 的临终标度，检验量纲，并说明灰体谱、物种阈值与Planck尺度为何使末期公式不能视作精确预测。
